%%%%%%%%%%%%%%%%%%%%%%%%%%%%%%%%%%%%%%%%%%%%%%%%%%%%%%%%%%%%%%%%%%%%%%
% Lớp 8 | Chương I. Đa thức | Phép cộng và phép trừ đa thức
% Dạng 8D1?3-2: Tìm đa thức thỏa đẳng thức cho trước
%%%%%%%%%%%%%%%%%%%%%%%%%%%%%%%%%%%%%%%%%%%%%%%%%%%%%%%%%%%%%%%%%%%%%%


% nguồn: De_AI_Toan8_cuoi_hoc_ki_i_1 (nhập 21/11/2026)
\begin{bt}%[8D1H3-2]%[ma:582f18d2]
	Tìm đa thức $A$ thỏa mãn $6x^2 - 3xy^2 + A = x^2 + y^2 - 2xy^2$, rồi tính giá trị của $A$ tại $x = 1$, $y = 2$.
	\loigiai{
		Ta có $A = (x^2 + y^2 - 2xy^2) - (6x^2 - 3xy^2) = x^2 + y^2 - 2xy^2 - 6x^2 + 3xy^2 = -5x^2 + xy^2 + y^2$. \\
Thay $x = 1$, $y = 2$ vào $A$, ta được $A = -5 \cdot 1^2 + 1 \cdot 2^2 + 2^2 = -5 + 4 + 4 = 3$.
	}
\end{bt}

% nguồn: De_AI_Toan8_cuoi_hoc_ki_i_1 (nhập 21/11/2026)
\begin{bt}%[8D1H3-2]%[ma:80392a2f]
	Tìm đa thức $B$ biết $B - (2xy - 4y^2) = 5xy + x^2 - 7y^2$, rồi tính giá trị của $B$ tại $x = 2$, $y = -1$.
	\loigiai{
		Ta có $B = (5xy + x^2 - 7y^2) + (2xy - 4y^2) = x^2 + 7xy - 11y^2$. \\
Thay $x = 2$, $y = -1$ vào $B$, ta được $B = 2^2 + 7 \cdot 2 \cdot (-1) - 11 \cdot (-1)^2 = 4 - 14 - 11 = -21$.
	}
\end{bt}

% nguồn: De_AI_Toan8_cuoi_hoc_ki_i_1 (nhập 21/11/2026)
\begin{bt}%[8D1H3-2]%[ma:51ef685f]
	Tìm đa thức $C$ biết $(4x^2 - 5xy + 3y^2) - C = 2x^2 + 3xy - y^2$.
	\loigiai{
		Ta có $C = (4x^2 - 5xy + 3y^2) - (2x^2 + 3xy - y^2)$ \\
$= 4x^2 - 5xy + 3y^2 - 2x^2 - 3xy + y^2 = 2x^2 - 8xy + 4y^2$.
	}
\end{bt}

% nguồn: De_AI_Toan8_cuoi_hoc_ki_i_1 (nhập 21/11/2026)
\begin{bt}%[8D1H3-2]%[ma:2802c7c0]
	Tìm hai đa thức $M$ và $N$ biết $M + N = 5x^2 - 2xy + y^2$ và $M - N = x^2 + 4xy - 3y^2$.
	\loigiai{
		Cộng vế theo vế hai đẳng thức: $2M = (5x^2 - 2xy + y^2) + (x^2 + 4xy - 3y^2) = 6x^2 + 2xy - 2y^2$. \\
Suy ra $M = 3x^2 + xy - y^2$. \\
Khi đó $N = (5x^2 - 2xy + y^2) - (3x^2 + xy - y^2) = 2x^2 - 3xy + 2y^2$.
	}
\end{bt}

% nguồn: De_AI_Toan8_cuoi_hoc_ki_i_1 (nhập 21/11/2026)
\begin{bt}%[8D1H3-2]%[ma:e27f961d]
	Tìm đa thức $A$ thỏa mãn $A \cdot 3xy = 6x^3y - 9x^2y^2 + 3xy^3$ (với $x \ne 0$, $y \ne 0$).
	\loigiai{
		Ta có $A = (6x^3y - 9x^2y^2 + 3xy^3) : (3xy) = 2x^2 - 3xy + y^2$.
	}
\end{bt}

% nguồn: De_AI_Toan8_cuoi_hoc_ki_i_1 (nhập 21/11/2026)
\begin{bt}%[8D1H3-2]%[ma:e12a92e5]
	Tìm đa thức $P$ biết $P : (-2x^2y) = 3x - 4y + 5$ (với $x \ne 0$, $y \ne 0$).
	\loigiai{
		Ta có $P = (-2x^2y)(3x - 4y + 5) = -6x^3y + 8x^2y^2 - 10x^2y$.
	}
\end{bt}

% nguồn: chuyen_dang (nhập 22/11/2026)
\begin{bt}%[8D1H3-2]%[ma:18c0055d]
	Cho  đa thức $P(x)=(1-x)^8$. Tính tổng các hệ số của đa thức $P(x)$.
	\loigiai{ 
		Ta có  $P(x)=(1-x)^8=\displaystyle\displaystyle\sum\limits_{k=0}^8 \mathrm{C}_8^k(-1)^k x^k$.\\
		Khi đó tổng các hệ số của đa thức  $P(x)$ là
		 $S=\mathrm{C}_8^0-\mathrm{C}_8^1+\cdots-\mathrm{C}_8^7+\mathrm{C}_8^8$.\\
		Xét khai triển $(a+b)^8=\displaystyle\sum\limits_{k=0}^8  \mathrm{C}_8^k a^{8-k} b^k$.\\
		Ta chọn  $a=1$, $b=-1$ thu được $(1-1)^8=\mathrm{C}_8^0-\mathrm{C}_8^1+\mathrm{C}_8^2-\cdots-\mathrm{C}_8^7+\mathrm{C}_8^8$.\\
		Vậy  tổng các hệ số của đa thức $P(x)$ bằng $0$.
	 }
\end{bt}

% nguồn: chuyen_dang (nhập 22/11/2026)
\begin{bt}%[8D1H3-2]%[ma:789bb18d]
	Tìm đa thức $M$ biết $M + (x - 2)(x + 3) = 2x^2 - 5x + 1$.
	\loigiai{
		Ta có $(x - 2)(x + 3) = x^2 + x - 6$. \\
Suy ra $M = (2x^2 - 5x + 1) - (x^2 + x - 6) = 2x^2 - 5x + 1 - x^2 - x + 6 = x^2 - 6x + 7$.
	}
\end{bt}
